\chapter{Αποτελέσματα Αξιολόγησης}
\label{ch:experiments}

Τα αποτελέσματα παρουσιάζονται σε δύο μέρη. Πρώτα η στατική εκπαίδευση, με ημερομηνία
αποκοπής στις 30 Νοεμβρίου 2025 και περίοδο αξιολόγησης τον Δεκέμβριο 2025. Η
προσέγγιση αυτή χρησιμεύει ως βάση αναφοράς για τη σύγκριση κάθε στρατηγικής κάτω από
σταθερές συνθήκες. Έπειτα η μηνιαία walk-forward επανεκπαίδευση στο Q1 2026, η οποία
επιτρέπει στα μοντέλα να προσαρμόζονται στις μεταβαλλόμενες συνθήκες και διακυμάνσεις
της αγοράς.

\section{Ρύθμιση Πειραμάτων}
\label{sec:exp_setup}

Στην Αγορά Επόμενης Ημέρας [Day-Ahead Market (DAM)], οι προσφορές κλείνουν στις 12:00 CET
της προηγούμενης ημέρας (D-1). Η στρατηγική Teacher-Forcing (TF) χρησιμοποιεί τις πραγματικές
τιμές σε κάθε βήμα και λειτουργεί ως μοντέλο αναφοράς (oracle) που ορίζει το θεωρητικό άνω
φράγμα της απόδοσης. Οι στρατηγικές Recursive, MIMO και Direct προβλέπουν ολόκληρο το 24ωρο
χρησιμοποιώντας μόνο τη διαθέσιμη πληροφορία πριν από το κλείσιμο της πύλης (gate closure).

Η ανάλυση σημαντικότητας χαρακτηριστικών [feature importance] (Κεφ.~\ref{ch:fi})
επιβεβαιώνει αυτή τη διάκριση: στην Teacher-Forcing η πραγματική τιμή της προηγούμενης ώρας
συγκεντρώνει το 88\% της βαρύτητας, ενώ οι διαθέσιμες προβλέψεις στο κλείσιμο της πύλης
(φορτίο και ΑΠΕ επόμενης ημέρας) επηρεάζουν αθροιστικά μόλις κατά 0,6\% (βλ.\ Κεφ.~\ref{ch:fi}). Η υψηλή επίδοση
της TF οφείλεται αποκλειστικά στην πρόσβαση στα πραγματικά προηγούμενα δεδομένα, επομένως
δεν αποτελεί ρεαλιστική εκτίμηση για τις πραγματικές συνθήκες της αγοράς.

Τα μοντέλα Μηχανικής Μάθησης εκπαιδεύονται από τον Ιανουάριο 2017 έως τις 30 Νοεμβρίου 2025,
με βελτιστοποίηση υπερπαραμέτρων μέσω Optuna. Εξαίρεση αποτελεί το μοντέλο SVR, το οποίο
σε ορισμένες περιπτώσεις εκπαιδεύεται σε περιορισμένο εύρος 5.000 ωρών για τη μείωση του
υπολογιστικού κόστους. Η επίδοση συγκρίνεται με τέσσερα βασικά μοντέλα αναφοράς (Naive-1,
Naive-24, Naive-168 και Seasonal Profile). Ο Πίνακας~\ref{tab:benchmarks_performance}
συγκεντρώνει τις τιμές τους για τις δύο περιόδους ελέγχου (Δεκέμβριος 2025: 744 ώρες·
Q1 2026: 2.160 ώρες). Με τον όρο Q1 2026 αναφερόμαστε στο τρίμηνο δοκιμής
Δεκέμβριος 2025 - Φεβρουάριος 2026, και όχι στο ημερολογιακό πρώτο τρίμηνο
(Ιανουάριος - Μάρτιος).

\begin{table}[H]
  \centering
  \caption{Επίδοση μοντέλων αναφοράς (benchmarks) για την τιμή DAM και το ηλεκτρικό
           φορτίο, μηνιαία αξιολόγηση.}
  \label{tab:benchmarks_performance}
  \small
  \begin{tabular}{llcccc}
    \toprule
    \textbf{Περίοδος} & \textbf{Στόχος} & \textbf{Μοντέλο Αναφοράς} & \textbf{MAE} & \textbf{RMSE} & \textbf{sMAPE (\%)} \\
    \midrule
    Δεκ.\ 2025 & Τιμή (€/MWh) & Naive-1 & 13,709 & 22,748 & 13,013 \\
     & & Seasonal Profile & 18,127 & 27,388 & 17,293 \\
     & & Naive-24 & 20,032 & 31,257 & 20,288 \\
     & & Naive-168 & 27,504 & 40,912 & 27,324 \\
    \cmidrule{2-6}
     & Φορτίο (MW) & Naive-1 & 260,077 & 328,760 & 4,679 \\
     & & Seasonal Profile & 291,557 & 396,534 & 5,038 \\
     & & Naive-24 & 309,920 & 435,517 & 5,430 \\
     & & Naive-168 & 313,240 & 419,529 & 5,417 \\
    \midrule
    Q1 2026 & Τιμή (€/MWh) & Naive-1 & 15,433 & 26,781 & 22,635 \\
     & & Naive-24 & 23,321 & 37,008 & 34,260 \\
     & & Seasonal Profile & 25,254 & 38,377 & 31,244 \\
     & & Naive-168 & 29,204 & 44,537 & 37,895 \\
    \cmidrule{2-6}
     & Φορτίο (MW) & Naive-1 & 271,185 & 344,041 & 4,761 \\
     & & Naive-24 & 345,694 & 487,585 & 5,909 \\
     & & Seasonal Profile & 348,502 & 494,348 & 5,832 \\
     & & Naive-168 & 369,785 & 530,274 & 6,229 \\
    \bottomrule
  \end{tabular}
\end{table}

\noindent Οι τιμές του Πίνακα~\ref{tab:benchmarks_performance} αφορούν πρόβλεψη ενός
βήματος (Teacher-Forcing, Recursive). Από τα μοντέλα αναφοράς, το Naive-1 δίνει το
χαμηλότερο σφάλμα (13,7~€/MWh, τιμή Δεκ.), καθώς η τιμή μιας ώρας μοιάζει πολύ με
εκείνη της επόμενης. Το Naive-24 (20,0~€/MWh) αντιγράφει την προηγούμενη ημέρα και
αστοχεί όταν η τιμή ανεβαίνει ή πέφτει από μέρα σε μέρα, ενώ το Naive-168
(27,5~€/MWh) έχει τη χειρότερη απόδοση, επαναλαμβάνοντας την αντίστοιχη ώρα της
προηγούμενης εβδομάδας. Το Seasonal Profile (18,1~€/MWh), ως μέσος όρος τεσσάρων
εβδομάδων, δίνει μια ομαλή αλλά χονδρική εκτίμηση. Στο Q1 όλες οι αναφορές
χειροτερεύουν για την τιμή, ενώ για το φορτίο μένουν σχεδόν αμετάβλητες, ένδειξη ότι
η δυσκολία της περιόδου αφορά κυρίως την τιμή.

\noindent Για τις στρατηγικές ταυτόχρονης ημερήσιας πρόβλεψης
(MIMO, Direct) η φυσική αναφορά είναι το Naive-24, δηλαδή η ίδια ώρα της προηγούμενης
ημέρας, υπολογισμένη στο αντίστοιχο ημερήσιο σχήμα $H=24$: 18,212~€/MWh για την
τιμή και 309,920~MW για το φορτίο τον Δεκέμβριο 2025. Σε αυτό το πλαίσιο το Naive-1
δεν αποτελεί ισοδύναμη αναφορά, καθώς προϋποθέτει πρόσβαση στην αμέσως προηγούμενη
πραγματική ώρα.

\noindent Η αξιολόγηση των στρατηγικών πραγματοποιείται σε δύο στάδια. Αρχικά, οι
αρχιτεκτονικές TF, Recursive, MIMO και Direct εξετάζονται υπό στατική εκπαίδευση με βάση
αποκλειστικά τα αποτελέσματα του Δεκεμβρίου 2025. Ο μήνας αυτός χρησιμοποιείται ως ελεγχόμενη
περίοδος αναφοράς (control period), επιτρέποντας τη σύγκριση της δομικής συμπεριφοράς των
αλγορίθμων υπό σταθερές συνθήκες. Στη συνέχεια, η ανάλυση μεταφέρεται στο τρίμηνο
Q1 2026. Λόγω της μεταβολής καθεστώτος (regime shift) κατά την περίοδο αυτή, οι στατικές
επιδόσεις δεν αναλύονται σε επίπεδο αναλυτικών πινάκων ανά στρατηγική· αντ' αυτού,
συμπυκνώνονται και αντιπαραβάλλονται απευθείας με τις επιδόσεις της μηνιαίας
walk-forward επανεκπαίδευσης, μέσω του δείκτη Retraining Premium και του τελικού
συγκεντρωτικού πίνακα.

% ---------------------------------------------------------------
\section{Στατική Εκπαίδευση: Δεκέμβριος 2025}
\label{sec:static_dec}

Η στατική αξιολόγηση του Δεκεμβρίου 2025 αποτελεί το πρώτο πειραματικό στάδιο του
Κεφαλαίου. Τα μοντέλα εκπαιδεύονται μία φορά, με ημερομηνία αποκοπής την 30ή
Νοεμβρίου 2025, και αξιολογούνται στον Δεκέμβριο 2025, δηλαδή σε 744 ωριαίες
παρατηρήσεις. Με τον τρόπο αυτό εξετάζεται η συμπεριφορά των στρατηγικών πρόβλεψης
υπό κοινό και σταθερό εκπαιδευτικό πλαίσιο, πριν εισαχθεί η μηνιαία walk-forward
επανεκπαίδευση του Q1 2026.

Η ενότητα οργανώνεται ανά στρατηγική πρόβλεψης. Πρώτα παρουσιάζεται το
Teacher-Forcing, το οποίο λειτουργεί ως ιδεατό μοντέλο αναφοράς. Στη συνέχεια
εξετάζεται η Recursive στρατηγική, που αποτελεί πρακτικά εφαρμόσιμη αναδρομική
προσέγγιση. Τέλος, παρουσιάζονται οι στρατηγικές MIMO και Direct, οι οποίες
προβλέπουν ολόκληρο το 24ωρο από μία κοινή ημερήσια αφετηρία. Στο τέλος της ενότητας
παρατίθεται συγκεντρωτική σύνοψη των καλύτερων αποτελεσμάτων κάθε στρατηγικής για
την τιμή DAM και το ηλεκτρικό φορτίο.

\subsection{Teacher-Forcing (TF)}
\label{subsec:tf_static_dec}

Η στρατηγική Teacher-Forcing παράγει τις ωριαίες προβλέψεις διαδοχικά,
χρησιμοποιώντας την πραγματική τιμή του προηγούμενου βήματος ως είσοδο. Με τον
τρόπο αυτό εξαλείφει τη συσσώρευση σφάλματος (error accumulation) και λειτουργεί ως
ιδεατό μοντέλο αναφοράς (oracle). Ωστόσο, απαιτεί τη σταδιακή διαθεσιμότητα των
πραγματικών δεδομένων, κάτι που δεν ισχύει στο gate closure της Αγοράς Επόμενης
Ημέρας. Επομένως, η επίδοση της TF δεν πρέπει να ερμηνεύεται ως πλήρως αναπτυξίσιμη
επιχειρησιακή λύση, αλλά ως άνω φράγμα απόδοσης.

Οι Πίνακες~\ref{tab:tf_price_dec} και~\ref{tab:tf_load_dec} παρουσιάζουν τα
αποτελέσματα Teacher-Forcing για τον Δεκέμβριο 2025.

\begin{table}[H]
  \centering
  \caption{TF, πρόβλεψη τιμής, Δεκέμβριος 2025 (744h).}
  \label{tab:tf_price_dec}
  \begin{tabular}{lcccc}
    \toprule
    \textbf{Μοντέλο} & \textbf{MAE (€/MWh)} & \textbf{RMSE (€)} & \textbf{sMAPE (\%)} & \textbf{Κατ.} \\
    \midrule
    Ensemble-Best3        & 9,943  & 15,994 & 10,697 & 1 \\
    XGBoost               & 10,204 & 16,170 & 11,087 & 2 \\
    LightGBM              & 10,334 & 16,236 & 11,178 & 3 \\
    RandomForest          & 11,220 & 18,206 & 11,467 & 4 \\
    MLP-Optuna            & 11,737 & 17,191 & 12,371 & 5 \\
    SVR                   & 13,432 & 19,070 & 13,666 & 6 \\
    \midrule
    Naive-1 (baseline)    & 13,709 & 22,748 & 13,013 & --- \\
    \bottomrule
  \end{tabular}
\end{table}

\begin{table}[H]
  \centering
  \caption{TF, πρόβλεψη φορτίου, Δεκέμβριος 2025 (744h).}
  \label{tab:tf_load_dec}
  \begin{tabular}{lcccc}
    \toprule
    \textbf{Μοντέλο} & \textbf{MAE (MW)} & \textbf{RMSE (MW)} & \textbf{sMAPE (\%)} & \textbf{Κατ.} \\
    \midrule
    Ensemble-Best3        & 72,916  & 96,945  & 1,260 & 1 \\
    XGBoost               & 75,885  & 98,951  & 1,311 & 2 \\
    LightGBM              & 76,324  & 101,752 & 1,316 & 3 \\
    MLP-Optuna            & 78,992  & 106,963 & 1,370 & 4 \\
    RandomForest          & 96,539  & 133,149 & 1,671 & 5 \\
    SVR                   & 141,406 & 199,114 & 2,469 & 6 \\
    \midrule
    Naive-1 (baseline)    & 260,077 & 328,760 & 4,679 & --- \\
    Seasonal Profile      & 291,557 & 396,534 & 5,038 & --- \\
    \bottomrule
  \end{tabular}
\end{table}

Στο πλαίσιο της στρατηγικής Teacher-Forcing, το Ensemble-Best3 επιτυγχάνει την καλύτερη
επίδοση τόσο στην τιμή όσο και στο φορτίο. Στην πρόβλεψη της τιμής DAM καταγράφει
MAE 9,943~€/MWh, ενώ τα XGBoost και LightGBM ακολουθούν με μικρή διαφορά. Το αποτέλεσμα
αυτό επιβεβαιώνει τη σταθερή συμπεριφορά των gradient boosting μοντέλων στο πρόβλημα
της τιμής, αλλά και την πρόσθετη αξία του συνδυασμού μοντέλων.

Στο ηλεκτρικό φορτίο, το Ensemble-Best3 επιτυγχάνει MAE 72,916~MW και sMAPE 1,260\%.
Το σχετικό σφάλμα είναι πολύ χαμηλότερο από εκείνο της τιμής, γεγονός που αντανακλά τη
μεγαλύτερη κανονικότητα, εποχικότητα και αυτοσυσχέτιση του φορτίου. Η SVR παρουσιάζει
τη χαμηλότερη επίδοση μεταξύ των μοντέλων μηχανικής μάθησης, ιδίως στο φορτίο,
καθώς εκπαιδεύτηκε σε περιορισμένο ιστορικό εύρος για τη μείωση του υπολογιστικού
κόστους.

\subsection{Recursive (Rec)}
\label{subsec:rec_static_dec}

Η Recursive στρατηγική χρησιμοποιεί ένα ενιαίο μοντέλο που προβλέπει μία χρονική στιγμή
τη φορά και στη συνέχεια εισάγει την προηγούμενη πρόβλεψη ως είσοδο για την επόμενη
ώρα. Η προσέγγιση αυτή είναι πρακτικά εφαρμόσιμη στον ορίζοντα της DAM, καθώς δεν
απαιτεί πρόσβαση σε πραγματικές ενδοημερήσιες τιμές της ημέρας πρόβλεψης. Το βασικό
της μειονέκτημα είναι η πιθανή συσσώρευση σφάλματος όσο η πρόβλεψη προχωρά βαθύτερα
στον 24ωρο ορίζοντα.

\begin{table}[H]
  \centering
  \caption{Rec $H=24$, πρόβλεψη τιμής, Δεκέμβριος 2025 (31 ημερήσια τμήματα).}
  \label{tab:rec_price_dec}
  \small
  \begin{tabular}{lcccc}
    \toprule
    \textbf{Μοντέλο} & \textbf{MAE (€/MWh)} & \textbf{RMSE (€)} & \textbf{sMAPE (\%)} & \textbf{Κατ.} \\
    \midrule
    Ensemble-Top3 (LGBM+RF+XGB)      & \textbf{13,960} & 21,958 & 13,649 & 1 \\
    Ensemble-All-24h (1/MAE)          & 13,996 & 21,641 & 13,771 & 2 \\
    LGBM-Daily-Optuna                 & 14,494 & 22,511 & 14,208 & 3 \\
    RF-Daily-SS                       & 14,607 & 23,015 & 14,072 & 4 \\
    XGB-Daily-SS-Optuna Dense         & 14,816 & 23,088 & 14,743 & 5 \\
    MLP-SS                            & 15,169 & 22,065 & 14,791 & 6 \\
    SVR-Rec                           & 18,549 & 25,602 & 17,608 & 7 \\
    \bottomrule
  \end{tabular}
\end{table}

\begin{table}[H]
  \centering
  \caption{Rec $H=24$, πρόβλεψη φορτίου, Δεκέμβριος 2025 (31 ημερήσια τμήματα).}
  \label{tab:rec_load_dec}
  \begin{tabular}{lcccc}
    \toprule
    \textbf{Μοντέλο} & \textbf{MAE (MW)} & \textbf{RMSE (MW)} & \textbf{sMAPE (\%)} & \textbf{Κατ.} \\
    \midrule
    LGBM-Daily-SS-Optuna  & \textbf{80,704}  & 106,076 & 1,387 & 1 \\
    Ensemble-Top3 (1/MAE) & 92,126  & 121,919 & 1,581 & 2 \\
    XGB-Daily-Optuna      & 99,575  & 131,606 & 1,702 & 3 \\
    RF-Daily-SS           & 117,740 & 164,388 & 2,034 & 4 \\
    MLP-SS                & 122,604 & 160,775 & 2,108 & 5 \\
    SVR-Rec               & 184,739 & 254,578 & 3,220 & 6 \\
    \bottomrule
  \end{tabular}
\end{table}

Στην πρόβλεψη της τιμής DAM, το Ensemble-Top3 επιτυγχάνει την καλύτερη αναδρομική
επίδοση με MAE 13,960~€/MWh. Η τιμή αυτή βρίσκεται πολύ κοντά στο Naive-1
(13,709~€/MWh), αλλά δεν το υπερβαίνει. Το αποτέλεσμα αυτό δείχνει ότι, στη στατική
εκπαίδευση του Δεκεμβρίου, η Recursive στρατηγική περιορίζεται από τη συσσώρευση
σφάλματος. Παρότι τα δενδρικά μοντέλα αξιοποιούν μη γραμμικές σχέσεις και εξωγενή
χαρακτηριστικά, η αναδρομική χρήση των ίδιων των προβλέψεών τους αυξάνει το σφάλμα
κατά την εξέλιξη του 24ωρου ορίζοντα.

Στο ηλεκτρικό φορτίο, αντίθετα, η Recursive στρατηγική εμφανίζει πολύ ισχυρότερη
συμπεριφορά. Το LGBM-Daily-SS-Optuna επιτυγχάνει MAE 80,704~MW και πλησιάζει αρκετά
την επίδοση του Teacher-Forcing. Η διαφορά αυτή εξηγείται από τη σταθερότερη ημερήσια
και εβδομαδιαία δομή του φορτίου, η οποία περιορίζει την αποσταθεροποίηση των
αναδρομικών εισόδων. Η Προγραμματισμένη Δειγματοληψία (Scheduled Sampling) συμβάλλει
επιπλέον στη μείωση του χάσματος μεταξύ συνθηκών εκπαίδευσης και συνθηκών εκτέλεσης.

\subsection{MIMO και Direct}
\label{subsec:mimo_direct_static_dec}

Οι στρατηγικές MIMO και Direct προβλέπουν ολόκληρο το 24ωρο από μία κοινή ημερήσια
αφετηρία, χωρίς αναδρομική ανατροφοδότηση με ενδιάμεσες προβλέψεις. Για τον λόγο αυτό
η καταλληλότερη αναφορά τους δεν είναι το Naive-1, αλλά το Naive-24, δηλαδή η τιμή της
ίδιας ώρας της προηγούμενης ημέρας. Η σύγκρισή τους με το Naive-24 είναι μεθοδολογικά
πιο δίκαιη, καθώς μοιράζεται το ίδιο βασικό σύνολο διαθέσιμης πληροφορίας.

\paragraph{MIMO.}
Η στρατηγική MIMO (Multi-Input Multi-Output) χρησιμοποιεί ένα ενιαίο μοντέλο για την
ταυτόχρονη παραγωγή του διανύσματος των 24 ωρών της επόμενης ημέρας. Με τον τρόπο
αυτό αποφεύγει τη συσσώρευση σφάλματος, επειδή οι προβλέψεις δεν τροφοδοτούνται
αναδρομικά στο μοντέλο. Ωστόσο, δεν διαθέτει ενδοημερήσια ανατροφοδότηση και άρα
βασίζεται αποκλειστικά στην πληροφορία που είναι διαθέσιμη στην αρχή του ημερήσιου
ορίζοντα.

\begin{table}[H]
  \centering
  \caption{MIMO $H=24$, πρόβλεψη τιμής, Δεκέμβριος 2025.}
  \label{tab:mimo_price_dec}
  \begin{tabular}{lcccc}
    \toprule
    \textbf{Μοντέλο} & \textbf{MAE (€/MWh)} & \textbf{RMSE (€)} & \textbf{sMAPE (\%)} & \textbf{Κατ.} \\
    \midrule
    Ensemble-Best3 (1/MAE) & \textbf{14,512} & 21,560 & 13,568 & 1 \\
    XGB MIMO Dense         & 14,919 & 22,044 & 13,852 & 2 \\
    RF MIMO+Load Dense     & 15,139 & 22,977 & 13,901 & 3 \\
    MLP MIMO+Load Dense    & 19,325 & 25,646 & 18,612 & 4 \\
    SVR MIMO+Load          & 20,410 & 30,598 & 18,936 & 5 \\
    \midrule
    Naive-24 (αναφορά)     & 18,212 & 27,436 & 17,741 & --- \\
    Seasonal Profile       & 18,392 & 27,458 & 17,333 & --- \\
    \bottomrule
  \end{tabular}
\end{table}

Στην πρόβλεψη της τιμής, το Ensemble-Best3 ηγείται με MAE 14,512~€/MWh, επιτυγχάνοντας
μείωση σφάλματος κατά 20,3\% έναντι του Naive-24. Οι tree-based dense υλοποιήσεις
ακολουθούν με μικρή απόσταση, ενώ τα μοντέλα MLP και SVR υστερούν έναντι της αναφοράς.

\begin{table}[H]
  \centering
  \caption{MIMO $H=24$, πρόβλεψη φορτίου, Δεκέμβριος 2025. Το $\dagger$
           δηλώνει μοντέλο που εξαιρείται λόγω αστοχίας.}
  \label{tab:mimo_load_dec}
  \begin{tabular}{lcccc}
    \toprule
    \textbf{Μοντέλο} & \textbf{MAE (MW)} & \textbf{RMSE (MW)} & \textbf{sMAPE (\%)} & \textbf{Κατ.} \\
    \midrule
    Ensemble-Best3 (1/MAE) & \textbf{188,078} & 294,694 & 3,218 & 1 \\
    RF MIMO                & 219,250 & 313,966 & 3,758 & 2 \\
    MLP MIMO               & 219,884 & 308,093 & 3,789 & 3 \\
    SVR MIMO               & 586,623 & 709,802 & 10,406 & $\dagger$ \\
    \midrule
    Seasonal Profile       & 291,557 & 396,534 & 5,038 & --- \\
    Naive-24 (αναφορά)     & 309,920 & 435,517 & 5,430 & --- \\
    \bottomrule
  \end{tabular}
\end{table}

Στο ηλεκτρικό φορτίο, το Ensemble-Best3 παραμένει κορυφαίο με MAE 188,078~MW,
σημειώνοντας βελτίωση 39,3\% έναντι του Naive-24. Η SVR καταρρέει
(586,623~MW, sMAPE 10,406\%), γεγονός που δείχνει τη μειωμένη ανθεκτικότητα του
RBF πυρήνα στο συγκεκριμένο σχήμα εκπαίδευσης.

\paragraph{Direct.}
Η Direct στρατηγική εκπαιδεύει 24 ανεξάρτητα υπομοντέλα, ένα για κάθε ώρα της επόμενης
ημέρας. Η προσέγγιση αυτή επιτρέπει εξειδικευμένη προσαρμογή στα ωριαία μοτίβα,
θυσιάζοντας όμως τη χρονική συνοχή μεταξύ διαδοχικών προβλέψεων.

\begin{table}[H]
  \centering
  \caption{Direct $H=24$, πρόβλεψη τιμής, Δεκέμβριος 2025.}
  \label{tab:direct_price_dec}
  \begin{tabular}{lcccc}
    \toprule
    \textbf{Μοντέλο} & \textbf{MAE (€/MWh)} & \textbf{RMSE (€)} & \textbf{sMAPE (\%)} & \textbf{Κατ.} \\
    \midrule
    LGBM Direct Dense        & \textbf{14,543} & 21,788 & 13,559 & 1 \\
    Ensemble-All-24h (1/MAE) & 14,641 & 21,499 & 13,683 & 2 \\
    XGB Direct Dense         & 15,092 & 21,698 & 14,125 & 3 \\
    \midrule
    Naive-24 (αναφορά)       & 18,212 & 27,436 & 17,741 & --- \\
    Seasonal Profile         & 18,392 & 27,458 & 17,333 & --- \\
    \bottomrule
  \end{tabular}
\end{table}

Στην πρόβλεψη της τιμής, το LGBM Direct Dense προηγείται με MAE 14,543~€/MWh,
επιτυγχάνοντας μείωση σφάλματος κατά 20,1\% έναντι του Naive-24.

\begin{table}[H]
  \centering
  \caption{Direct $H=24$, πρόβλεψη φορτίου, Δεκέμβριος 2025.}
  \label{tab:direct_load_dec}
  \begin{tabular}{lcccc}
    \toprule
    \textbf{Μοντέλο} & \textbf{MAE (MW)} & \textbf{RMSE (MW)} & \textbf{sMAPE (\%)} & \textbf{Κατ.} \\
    \midrule
    XGB Direct Dense         & \textbf{188,033} & 296,611 & 3,218 & 1 \\
    Ensemble-All-24h (1/MAE) & 189,114 & 298,794 & 3,232 & 2 \\
    LGBM Direct Dense        & 194,417 & 304,411 & 3,314 & 3 \\
    \midrule
    Seasonal Profile         & 291,557 & 396,534 & 5,038 & --- \\
    Naive-24 (αναφορά)       & 309,920 & 435,517 & 5,430 & --- \\
    \bottomrule
  \end{tabular}
\end{table}

Στο ηλεκτρικό φορτίο, το XGB Direct Dense καταγράφει την καλύτερη επίδοση της
στρατηγικής με MAE 188,033~MW, δηλαδή πρακτικά ισοδύναμη με την καλύτερη MIMO
επίδοση. Συνολικά, οι MIMO και Direct εμφανίζουν πολύ κοντινά επίπεδα σφάλματος και
στις δύο εργασίες. Η προσθήκη πυκνών υστερήσεων (dense lags), αν και σχεδιάστηκε
αρχικά για την Direct αρχιτεκτονική, ωφελεί και τη MIMO στρατηγική, επειδή παρέχει
λεπτομερέστερη περιγραφή της ενδοημερήσιας δομής χωρίς αναδρομική διάδοση σφάλματος.

\subsection{Συνολική Σύνοψη Στατικής Αξιολόγησης}
\label{subsec:static_dec_summary}

Οι προηγούμενες υποενότητες παρουσιάζουν αναλυτικά τα αποτελέσματα κάθε στρατηγικής.
Για να είναι πιο άμεση η σύγκριση, οι Πίνακες~\ref{tab:static_best_price_dec} και
\ref{tab:static_best_load_dec} συγκεντρώνουν το καλύτερο αποτέλεσμα κάθε στρατηγικής
για την τιμή DAM και το ηλεκτρικό φορτίο αντίστοιχα. Στις στρατηγικές TF και Recursive
η σύγκριση γίνεται ως προς το Naive-1, ενώ στις MIMO και Direct ως προς το Naive-24,
επειδή οι τελευταίες προβλέπουν ολόκληρο το ημερήσιο προφίλ από μία κοινή αφετηρία.

\begin{table}[H]
  \centering
  \caption{Στατική εκπαίδευση, Δεκέμβριος 2025: καλύτερο μοντέλο ανά στρατηγική για την τιμή DAM.}
  \label{tab:static_best_price_dec}
  \small
  \setlength{\tabcolsep}{4pt}
  \begin{tabular}{llcccc}
    \toprule
    \textbf{Στρατηγική} & \textbf{Καλύτερο μοντέλο} &
    \textbf{MAE} & \textbf{RMSE} & \textbf{sMAPE (\%)} &
    \textbf{$\Delta$ MAE} \\
    \midrule
    TF &
    Ensemble-Best3 &
    9,943 &
    15,994 &
    10,697 &
    +27,5\% έναντι Naive-1 \\
    Recursive &
    Ensemble-Top3 &
    13,960 &
    21,958 &
    13,649 &
    -1,8\% έναντι Naive-1 \\
    MIMO &
    Ensemble-Best3 &
    14,512 &
    21,560 &
    13,568 &
    +20,3\% έναντι Naive-24 \\
    Direct &
    LGBM Direct Dense &
    14,543 &
    21,788 &
    13,559 &
    +20,1\% έναντι Naive-24 \\
    \bottomrule
  \end{tabular}
\end{table}

\begin{table}[H]
  \centering
  \caption{Στατική εκπαίδευση, Δεκέμβριος 2025: καλύτερο μοντέλο ανά στρατηγική για το ηλεκτρικό φορτίο.}
  \label{tab:static_best_load_dec}
  \small
  \setlength{\tabcolsep}{4pt}
  \begin{tabular}{llcccc}
    \toprule
    \textbf{Στρατηγική} & \textbf{Καλύτερο μοντέλο} &
    \textbf{MAE} & \textbf{RMSE} & \textbf{sMAPE (\%)} &
    \textbf{$\Delta$ MAE} \\
    \midrule
    TF &
    Ensemble-Best3 &
    72,916 &
    96,945 &
    1,260 &
    +72,0\% έναντι Naive-1 \\
    Recursive &
    LGBM-Daily-SS-Optuna &
    80,704 &
    106,076 &
    1,387 &
    +69,0\% έναντι Naive-1 \\
    MIMO &
    Ensemble-Best3 &
    188,078 &
    294,694 &
    3,218 &
    +39,3\% έναντι Naive-24 \\
    Direct &
    XGB Direct Dense &
    188,033 &
    296,611 &
    3,218 &
    +39,3\% έναντι Naive-24 \\
    \bottomrule
  \end{tabular}
\end{table}

\begin{figure}[H]
  \centering
  \includegraphics[width=\textwidth]{ch6_forecast_analysis/static_dec_best_price_mae.png}
  \caption{Στατική εκπαίδευση, Δεκέμβριος 2025: καλύτερο MAE ανά στρατηγική για την τιμή DAM.}
  \label{fig:static_dec_best_price_mae}
\end{figure}

\begin{figure}[H]
  \centering
  \includegraphics[width=\textwidth]{ch6_forecast_analysis/static_dec_best_load_mae.png}
  \caption{Στατική εκπαίδευση, Δεκέμβριος 2025: καλύτερο MAE ανά στρατηγική για το ηλεκτρικό φορτίο.}
  \label{fig:static_dec_best_load_mae}
\end{figure}

Η συνολική σύνοψη αναδεικνύει δύο βασικές διαφοροποιήσεις. Στην τιμή DAM, το
Teacher-Forcing έχει καθαρά την καλύτερη επίδοση, αλλά η υπεροχή του πρέπει να
ερμηνευθεί ως ιδεατό άνω φράγμα, επειδή χρησιμοποιεί πραγματικές ενδοημερήσιες
υστερήσεις. Από τις πρακτικά εφαρμόσιμες στρατηγικές, οι MIMO και Direct βελτιώνουν
αισθητά το Naive-24, ενώ η Recursive υπολείπεται οριακά του Naive-1. Αυτό δείχνει ότι
στην τιμή η αναδρομική διάδοση σφάλματος είναι ουσιαστικότερος περιορισμός από ό,τι
στο φορτίο.

Στο ηλεκτρικό φορτίο, η εικόνα είναι διαφορετική. Η Recursive στρατηγική πλησιάζει
αρκετά το Teacher-Forcing και υπερβαίνει κατά πολύ το Naive-1, γεγονός που δείχνει ότι
η ισχυρή ημερήσια και εβδομαδιαία κανονικότητα του φορτίου περιορίζει τη συσσώρευση
σφάλματος. Οι MIMO και Direct βελτιώνουν επίσης σημαντικά το Naive-24, όμως παραμένουν
αρκετά μακριά από τις στρατηγικές που αξιοποιούν διαδοχική ενδοημερήσια πληροφορία.
Τα Σχήματα~\ref{fig:static_dec_best_price_mae} και~\ref{fig:static_dec_best_load_mae}
οπτικοποιούν αυτή τη διαφορά: στην τιμή η απόσταση του Teacher-Forcing από τις
εφαρμόσιμες στρατηγικές είναι εμφανής, ενώ στο φορτίο η Recursive στρατηγική παραμένει
πολύ κοντά στο ιδεατό όριο. Η στατική αξιολόγηση, επομένως, δείχνει ήδη ότι η επιλογή
στρατηγικής είναι εξίσου σημαντική με την επιλογή αλγορίθμου και ότι τα δύο σήματα,
τιμή και φορτίο, δεν αντιδρούν με τον ίδιο τρόπο στον 24ωρο ορίζοντα πρόβλεψης.


% ---------------------------------------------------------------
\section{Walk-Forward Μηνιαία Επανεκπαίδευση, Q1 2026}
\label{sec:walkforward_q1}

Η walk-forward μηνιαία επανεκπαίδευση εφαρμόζεται κατά το πρώτο τρίμηνο του 2026 (Q1 2026: 2.160 ώρες αξιολόγησης), με σκοπό την ενσωμάτωση της πιο πρόσφατης πληροφορίας της αγοράς στα μοντέλα πρόβλεψης. Η διαδικασία αυτή επεκτείνει το σύνολο εκπαίδευσης (training set) κατά έναν πλήρη μήνα σε κάθε βήμα, επιτρέποντας στους αλγορίθμους να προσαρμόζουν τις παραμέτρους τους στις μεταβαλλόμενες συνθήκες. Στις υποενότητες που ακολουθούν, παρουσιάζονται αρχικά τα αποτελέσματα της στρατηγικής Teacher-Forcing (TF) υπό το πλαίσιο της μηνιαίας ανανέωσης δεδομένων, προκειμένου να ποσοτικοποιηθεί το θεωρητικό όφελος της δυναμικής εκπαίδευσης σε σχέση με τη στατική προσέγγιση.

\subsection{Αποτελέσματα Στρατηγικής Teacher-Forcing (TF) Walk-Forward}
\label{subsec:tf_wf_results}

Οι Πίνακες~\ref{tab:wf_tf_price} και~\ref{tab:wf_tf_load} συγκεντρώνουν τις συνολικές επιδόσεις των μοντέλων Μηχανικής Μάθησης για την τιμή DAM και το ηλεκτρικό φορτίο αντίστοιχα, κατά την περίοδο Q1 2026, χρησιμοποιώντας τη στρατηγική Teacher-Forcing με μηνιαία επανεκπαίδευση.

\begin{table}[H]
  \centering
  \caption{Walk-forward TF, πρόβλεψη τιμής DAM, Q1 2026 (2.160h).}
  \label{tab:wf_tf_price}
  \begin{tabular}{llccc}
    \toprule
    \textbf{Κατ.} & \textbf{Μοντέλο} & \textbf{MAE (€/MWh)} & \textbf{RMSE (€)} & \textbf{sMAPE (\%)} \\
    \midrule
    1   & TF-LGBM-MR          & \textbf{10,463} & 15,682 & 20,600 \\
    2   & Ensemble-TF-Best3   & 10,491          & 15,726 & 20,589 \\
    3   & TF-XGB-MR           & 10,913          & 16,148 & 21,008 \\
    4   & Ensemble-TF-All     & 11,041          & 16,214 & 20,901 \\
    5   & TF-RF-MR            & 12,806          & 19,089 & 21,825 \\
    6   & TF-MLP-MR           & 15,620          & 22,458 & 25,047 \\
    \midrule
    --- & Στατικό TF Best3    & 14,478          & 22,193 & 25,085 \\
    --- & Naive-1 (baseline)  & 15,433          & 26,781 & 22,635 \\
    \bottomrule
  \end{tabular}
\end{table}

Στην πρόβλεψη της τιμής DAM για το πρώτο τρίμηνο του 2026 (Πίνακας~\ref{tab:wf_tf_price}), το μοντέλο TF-LGBM-MR σημειώνει το χαμηλότερο MAE (10,463~€/MWh), με το Ensemble-TF-Best3 να ακολουθεί στα 10,491~€/MWh εξασφαλίζοντας παράλληλα το βέλτιστο sMAPE (20,589\%). Η σύγκριση αυτή με το Στατικό TF Best3 (14,478~€/MWh) μεταφράζεται σε μείωση του σφάλματος κατά 27,7\%, γεγονός που αναδεικνύει τη σημασία της επανεκπαίδευσης σε πρόσφατα δεδομένα.

Επειδή το κεντρικό αυτό αποτέλεσμα αφορά ακριβώς την τιμή, το σήμα όπου το \textit{gas\_price} της δοκιμαστικής περιόδου εμφανίζει ήπιο look-ahead λόγω της συνθετικής παρεμβολής (§\ref{ch:data}), σημειώνεται ότι το εύρημα δεν στηρίζεται σε αυτό το χαρακτηριστικό: στα νικητήρια μοντέλα TF/Recursive κυριαρχεί το αυτοπαλίνδρομο σήμα (\textit{Price\_lag1} έως $88{,}3\%$), ενώ το \textit{gas\_price} αποκτά μη αμελητέο βάρος μόνο στις πολυβηματικές δομές MIMO/Direct (Κεφ.~\ref{ch:fi}). Η επανεκπαίδευση κερδίζει μέσω της προσαρμογής στη μετατόπιση επιπέδου της τιμής (§\ref{subsec:regime_evidence}), όχι μέσω του αερίου· ο περιορισμός και η αντιμετώπισή του αναλύονται στο Κεφ.~\ref{ch:conclusions}.

\begin{table}[H]
  \centering
  \caption{Walk-forward TF, πρόβλεψη ηλεκτρικού φορτίου, Q1 2026 (2.160h).}
  \label{tab:wf_tf_load}
  \begin{tabular}{llccc}
    \toprule
    \textbf{Κατ.} & \textbf{Μοντέλο} & \textbf{MAE (MW)} & \textbf{RMSE (MW)} & \textbf{sMAPE (\%)} \\
    \midrule
    1   & Ensemble-TF-Best3  & \textbf{67,729} & 89,324  & 1,156 \\
    2   & TF-LGBM-MR         & 69,435          & 92,490  & 1,182 \\
    3   & TF-XGB-MR          & 72,925          & 95,540  & 1,242 \\
    4   & TF-MLP-MR          & 88,493          & 116,316 & 1,518 \\
    \midrule
    --- & Στατικό TF Best3   & 96,751          & 122,049 & 1,685 \\
    --- & Naive-1 (baseline) & 267,151         & 335,879 & 4,696 \\
    \bottomrule
  \end{tabular}
\end{table}

Στην πρόβλεψη του ηλεκτρικού φορτίου (Πίνακας~\ref{tab:wf_tf_load}), το Ensemble-TF-Best3 καταγράφει την καλύτερη επίδοση με MAE 67,729~MW και sMAPE 1,156\%. Η μηνιαία επανεκπαίδευση μειώνει το σφάλμα κατά 30\% σε σχέση με το Στατικό TF Best3 (96,751~MW).

Παρά τη μείωση του MAE, η ποσοστιαία μετρική sMAPE δείχνει ότι η πραγματική βελτίωση στο φορτίο είναι οριακή (μόλις $\Delta\text{sMAPE} = -0,529\%$). Αντίθετα, στην τιμή η πτώση του σφάλματος είναι πολύ μεγαλύτερη ($\Delta\text{sMAPE} = -4,496\%$). Η σύγκριση αυτή, επειδή γίνεται σε κοινή ποσοστιαία κλίμακα, εξαλείφει τη διαφορά των μονάδων μέτρησης (€/MWh έναντι MW) και αποδεικνύει ότι η επανεκπαίδευση ωφελεί ουσιαστικά την τιμή, κάτι που θα δειχθεί και στα σχήματα που ακολουθούν.

% ---------------------------------------------------------------
\subsection{Ανάλυση Μεταβολής Καθεστώτος}
\label{subsec:regime_evidence}

Η ασύμμετρη αυτή συμπεριφορά ανάμεσα στα δύο σήματα συνδέεται άμεσα με τη στατιστική εξέλιξη των πραγματικών δεδομένων. Η αιτία εντοπίζεται στην εξέλιξη της τιμής DAM και του ηλεκτρικού φορτίου κατά το πρώτο τρίμηνο του 2026, όπως αποτυπώνεται στα Σχήματα~\ref{fig:regime_changepoint} και~\ref{fig:regime_load}. Η μαύρη έντονη γραμμή αντιπροσωπεύει τον κυλιόμενο μέσο όρο 30 ημερών (level) και η πράσινη γραμμή την κυλιόμενη τυπική απόκλιση 30 ημερών, η οποία αποδίδει τη μεταβλητότητα (volatility) της αγοράς.

\begin{figure}[H]
  \centering
  \includegraphics[width=\textwidth]{ch6_forecast_analysis/30_regime_changepoint_price.png}
  \caption{Μεταβολή καθεστώτος Q1 2026: κυλιόμενος μέσος και κυλιόμενη μεταβλητότητα 30 ημερών του πραγματικού σήματος, Τιμή DAM (€/MWh).}
  \label{fig:regime_changepoint}
\end{figure}

Στο Σχήμα~\ref{fig:regime_changepoint} για την τιμή DAM, η γκρίζα ζώνη ορίζει τον Φεβρουάριο ως την περίοδο χαμηλών τιμών, ενώ η κόκκινη διακεκομμένη γραμμή στα μέσα του μήνα εντοπίζει στατιστικά το ακριβές σημείο της δομικής αλλαγής (change-point). Εδώ παρατηρείται μια βίαιη μεταβολή καθεστώτος (regime shift), καθώς η αγορά υποχωρεί από τις υψηλές τιμές της περιόδου εκπαίδευσης ($110 \pm 34$~€/MWh) και σταθεροποιείται στο νέο καθεστώς του Φεβρουαρίου στα 78~€/MWh (πτώση 28\%), με τη μεταβλητότητα όμως να αυξάνεται απότομα στα $\pm$45~€/MWh. Οι δύο αυτές κατανομές δεν επικαλύπτονται. Αυτό το covariate shift αποπροσανατολίζει τα στατικά μοντέλα. Έχοντας εκπαιδευτεί αποκλειστικά σε υψηλές τιμές, αδυνατούν να αναγνωρίσουν το νέο επίπεδο της αγοράς, κάνοντας τη walk-forward επανεκπαίδευση απολύτως αναγκαία.

\begin{figure}[H]
  \centering
  \includegraphics[width=\textwidth]{ch6_forecast_analysis/30_regime_changepoint_load.png}
  \caption{Μεταβολή καθεστώτος Q1 2026: κυλιόμενος μέσος και κυλιόμενη μεταβλητότητα 30 ημερών του πραγματικού σήματος, Ηλεκτρικό φορτίο (MW).}
  \label{fig:regime_load}
\end{figure}

Αντίθετα, στο Σχήμα~\ref{fig:regime_load} για το ηλεκτρικό φορτίο η εικόνα είναι διαφορετική, καθώς δεν εντοπίζεται κανένα στατιστικό σημείο αλλαγής. Αν και ο κυλιόμενος μέσος (level) σχηματίζει μια καμπύλη λόγω της φυσικής μετάβασης από τον χειμώνα στην άνοιξη, η μέση ζήτηση μεταβάλλεται ελάχιστα (πτώση μόλις 7\%), ενώ η μεταβλητότητα μειώνεται ελαφρώς. Το κρίσιμο εύρημα είναι ότι η ζώνη της εκπαίδευσης ($5707 \pm 951$~MW) και η ζώνη του Φεβρουαρίου ($5567 \pm 839$~MW) επικαλύπτονται σχεδόν ολοκληρωτικά. Το σήμα δεν βγαίνει ποτέ εκτός των ορίων που ήδη γνωρίζουν τα μοντέλα. Επειδή οι συνθήκες δεν άλλαξαν δομικά, τα στατικά μοντέλα παραμένουν «διαβασμένα», γεγονός που εξηγεί γιατί συνεχίζουν να προβλέπουν με υψηλή ακρίβεια και το όφελος της επανεκπαίδευσης είναι σχεδόν μηδενικό.

Για την αξιολόγηση των μοντέλων σε πραγματικές συνθήκες, η ανάλυση επεκτείνεται παρακάτω στις εφαρμόσιμες στρατηγικές πολλαπλών βημάτων (Recursive, MIMO, Direct) με μηνιαία επανεκπαίδευση, εξετάζοντας το Retraining Premium και τη συμπεριφορά τους στη μεταβολή καθεστώτος.

% ---------------------------------------------------------------
\subsection{Αντίκτυπος Επανεκπαίδευσης}
\label{subsec:retraining_impact}

Η σύγκριση των στατικών μοντέλων με τις walk-forward εκδόσεις τους ποσοτικοποιεί την επίδραση της μηνιαίας ανανέωσης των δεδομένων. Η ανάλυση εξετάζει την εξέλιξη του σφάλματος ανά μήνα για να προσδιορίσει τη στιγμή διαφοροποίησης των δύο προσεγγίσεων και στη συνέχεια μετρά τη συνολική ποσοστιαία μεταβολή των δεικτών MAE και sMAPE.

Ο Πίνακας~\ref{tab:wf_all_strategies} συγκεντρώνει το καλύτερο μοντέλο κάθε στρατηγικής υπό μηνιαία επανεκπαίδευση, για τιμή DAM και φορτίο, ώστε οι τέσσερις αρχιτεκτονικές να συγκρίνονται σε κοινή βάση. Η Teacher-Forcing παρατίθεται ως oracle αναφορά: χρησιμοποιεί τις πραγματικές πρόσφατες υστερήσεις και δεν είναι αναπτυξίσιμη στο κλείσιμο της πύλης. Η αναπτυξίσιμη σύσταση είναι η Recursive, η οποία με το \texttt{Rec-XGB-SS-DO-MR} επιτυγχάνει 18,534~€/MWh στην τιμή και 70,641~MW στο φορτίο, υπερέχοντας τόσο των πολυβηματικών στρατηγικών (MIMO, Direct) όσο και του Naive-1 και στις δύο εργασίες.

\begin{table}[H]
  \centering
  \caption{Walk-forward μηνιαία επανεκπαίδευση, Q1 2026 (2.160h): καλύτερο μοντέλο ανά στρατηγική για τιμή DAM και φορτίο.}
  \label{tab:wf_all_strategies}
  \begin{tabular}{llcccc}
    \toprule
    & & \multicolumn{2}{c}{\textbf{Τιμή DAM}} & \multicolumn{2}{c}{\textbf{Φορτίο}} \\
    \cmidrule(lr){3-4}\cmidrule(lr){5-6}
    \textbf{Στρατηγική} & \textbf{Μοντέλο} & \textbf{MAE} & \textbf{sMAPE} & \textbf{MAE} & \textbf{sMAPE} \\
                        &                  & (€/MWh) & (\%) & (MW) & (\%) \\
    \midrule
    Teacher-Forcing$^{\dagger}$        & TF-LGBM / Ens-Best3   & 10,463          & 20,600 & 67,729           & 1,156 \\
    Recursive$^{\ddagger}$             & Rec-XGB-SS-DO-MR      & \textbf{18,534} & 25,749 & \textbf{70,641}  & 1,204 \\
    MIMO                               & Ensemble-Best3 MR     & 19,012          & 26,423 & 206,786          & 3,456 \\
    Direct                             & LGBM Direct Dense MR  & 19,875          & 27,315 & 218,623          & 3,596 \\
    \midrule
    Naive-1 (baseline)                 & ---                   & 15,433          & 22,635 & 267,151          & 4,696 \\
    \bottomrule
  \end{tabular}

  \smallskip
  {\footnotesize $\dagger$: oracle, μη αναπτυξίσιμη (χρησιμοποιεί πραγματικές πρόσφατες υστερήσεις).\quad
  $\ddagger$: συνιστώμενη αναπτυξίσιμη στρατηγική.}
\end{table}

\begin{figure}[H]
  \centering
  \includegraphics[width=\textwidth]{ch6_forecast_analysis/11_monthly_price.png}
  \caption{Μηνιαία ανάλυση MAE στην τιμή DAM: Στατική εκπαίδευση έναντι Walk-Forward Monthly Retrain, Q1 2026.}
  \label{fig:monthly_price}
\end{figure}

Το Σχήμα~\ref{fig:monthly_price} δείχνει το MAE της τιμής DAM ανά μήνα, συγκρίνοντας τα στατικά μοντέλα (συμπαγείς μπάρες) με τα αντίστοιχα walk-forward (γραμμωτές μπάρες). Η διαφοροποίηση εμφανίζεται αμέσως τον Ιανουάριο, μετά τον πρώτο κύκλο επανεκπαίδευσης, όπου η walk-forward στρατηγική μειώνει άμεσα το σφάλμα. Το πλεονέκτημα αυτό διευρύνεται έντονα τον Φεβρουάριο κατά τη μεταβολή καθεστώτος, όπου η στατική Teacher-Forcing καταρρέει, ενώ η walk-forward συγκρατεί το σφάλμα. Η επανεκπαίδευση προσδίδει ανθεκτικότητα από τον πρώτο μήνα, η οποία εδραιώνεται με την πάροδο του χρόνου.

\begin{figure}[H]
  \centering
  \includegraphics[width=\textwidth]{ch6_forecast_analysis/12_monthly_load.png}
  \caption{Μηνιαία ανάλυση MAE στο ηλεκτρικό φορτίο: Στατική εκπαίδευση έναντι Walk-Forward Monthly Retrain, Q1 2026.}
  \label{fig:monthly_load}
\end{figure}

Το Σχήμα~\ref{fig:monthly_load} παρουσιάζει την αντίστοιχη μηνιαία πορεία για το ηλεκτρικό φορτίο, όπου η εικόνα αντιστρέφεται. Σε όλο το τρίμηνο, οι μπάρες της στατικής και της walk-forward εκπαίδευσης για τις ίδιες στρατηγικές έχουν σχεδόν ίδιο ύψος. Η επανεκπαίδευση δεν προσφέρει καμία διαφοροποίηση από μήνα σε μήνα. Το σταθερό στατιστικό καθεστώς στη ζήτηση καθιστά τη στατική εκπαίδευση εξίσου ισχυρή.

\begin{figure}[H]
  \centering
  \includegraphics[width=\textwidth]{ch6_forecast_analysis/25_cross_regime_delta_price.png}
  \caption{Συνολική διαφορά επίδοσης (Walk-Forward vs Static) ανά στρατηγική στην τιμή DAM, Q1 2026.}
  \label{fig:delta_price}
\end{figure}

Το Σχήμα~\ref{fig:delta_price} δείχνει τη συνολική μεταβολή του σφάλματος με τη walk-forward εκπαίδευση. Στην τιμή DAM, η επανεκπαίδευση έχει θετικό πρόσημο, καθώς καταγράφει θετικό αντίκτυπο στις τρεις από τις τέσσερις στρατηγικές, με τη MIMO και την Teacher-Forcing να παρουσιάζουν τη μεγαλύτερη βελτίωση. Αντίθετα, η Recursive είναι η μόνη που εμφανίζει αρνητική μεταβολή, η οποία όμως κρίνεται αμελητέα αφού η επίδραση στο συνολικό sMAPE περιορίζεται μόλις στο $-0,11$ ποσοστιαίες μονάδες (pp).

\begin{figure}[H]
  \centering
  \includegraphics[width=\textwidth]{ch6_forecast_analysis/26_cross_regime_delta_load.png}
  \caption{Συνολική διαφορά επίδοσης (Walk-Forward vs Static) ανά στρατηγική στο ηλεκτρικό φορτίο, Q1 2026.}
  \label{fig:delta_load}
\end{figure}

Στο Σχήμα~\ref{fig:delta_load} για το ηλεκτρικό φορτίο, ο αντίκτυπος της επανεκπαίδευσης είναι αμελητέος σε όλες τις στρατηγικές, είτε με θετικό είτε με αρνητικό πρόσημο. Η Teacher-Forcing και η Direct καταγράφουν μια οριακή αρνητική μεταβολή, ενώ η Recursive και η MIMO σημειώνουν μια εξίσου μικρή βελτίωση. Χαρακτηριστικό της σταθερότητας του σήματος είναι ότι η μεγαλύτερη μεταβολή στο σύνολο του φορτίου εντοπίζεται στην Direct και είναι μόλις $-0,15$~pp στον δείκτη sMAPE.

Συνολικά, η μηνιαία επανεκπαίδευση κρίνεται απαραίτητη και αποδοτική στην τιμή DAM, καθώς προστατεύει τα μοντέλα από δομικές μεταβολές καθεστώτος. Αντίθετα, στο ηλεκτρικό φορτίο ο αντίκτυπος είναι αμελητέος, αφού η σταθερότητα του σήματος επιτρέπει στα στατικά μοντέλα να διατηρούν την ακρίβειά τους.

% ---------------------------------------------------------------
\section{Συγκριτική Αξιολόγηση Στρατηγικών Πρόβλεψης}
\label{sec:strategic_evaluation}

Η αξιολόγηση μετατοπίζεται πλέον από τους συγκεντρωτικούς δείκτες σφάλματος στην εις βάθος εξέταση της συμπεριφοράς των τεσσάρων στρατηγικών πρόβλεψης. Οι μέσες μετρικές MAE, RMSE και sMAPE αποτυπώνουν τη συνολική απόδοση, αδυνατούν όμως να αποκαλύψουν πού και πώς αστοχεί κάθε αρχιτεκτονική μέσα στο 24ωρο ή στις κρίσιμες στιγμές της αγοράς. Για τον λόγο αυτό, η ανάλυση δομείται σε τρεις άξονες: τη στατιστική φύση και μεροληψία των σφαλμάτων, τη χρονική εξέλιξη της ακρίβειας σε ωριαίο και κυκλικό επίπεδο, και τέλος την ανθεκτικότητα των αρχιτεκτονικών στις ακραίες συνθήκες αιχμής της αγοράς.

\subsection{Στατιστική Κατανομή και Μεροληψία Σφαλμάτων}
\label{subsubsect:statistical_bias}

Τα διαγράμματα πλαισίου-απολήξεων (box plots) αποτυπώνουν γραφικά τη στατιστική κατανομή των απόλυτων σφαλμάτων πρόβλεψης. Ένα στενό κεντρικό πλαίσιο με κοντές απολήξεις υποδηλώνει μια σταθερή και αξιόπιστη στρατηγική. Αυτό συμβαίνει διότι το πλαίσιο ορίζει το ενδοτεταρτημοριακό εύρος (IQR) όπου συγκεντρώνεται το 50\% των σφαλμάτων, η εσωτερική γραμμή δείχνει τη διάμεσο (median), ενώ οι οριζόντιες απολήξεις (whiskers) περιορίζουν το εύρος των μεγάλων σφαλμάτων από την 5η έως την 95η εκατοστιαία θέση.

\begin{figure}[H]
  \centering
  \includegraphics[width=\textwidth]{ch6_forecast_analysis/05_boxplot_price.png}
  \caption{Κατανομή απόλυτων σφαλμάτων πρόβλεψης στην τιμή DAM, Q1 2026.}
  \label{fig:boxplot_price}
\end{figure}

Το Σχήμα~\ref{fig:boxplot_price} δείχνει την κατανομή των σφαλμάτων για την τιμή DAM. Οι μονοβηματικές στρατηγικές (Teacher-Forcing και Recursive) επικρατούν με στενότερα πλαίσια και χαμηλότερα σφάλματα. Ωστόσο, ακόμη και σε αυτές, οι απολήξεις (whiskers) εκτείνονται αναλογικά πολύ ψηλότερα από τη διάμεσο. Το γεγονός αυτό φανερώνει ότι στην αγορά της τιμής DAM τα ακραία σφάλματα παραμένουν πολύ έντονα γενικά, λόγω των ξαφνικών εξάρσεων (spikes). Αντίθετα, οι πολυβηματικές προσεγγίσεις (MIMO και Direct) εμφανίζουν πολύ πιο απλωμένες κατανομές, επιβεβαιώνοντας τη μειωμένη σταθερότητά τους.

\begin{figure}[H]
  \centering
  \includegraphics[width=\textwidth]{ch6_forecast_analysis/06_boxplot_load.png}
  \caption{Κατανομή απόλυτων σφαλμάτων πρόβλεψης στο ηλεκτρικό φορτίο, Q1 2026.}
  \label{fig:boxplot_load}
\end{figure}

Στο Σχήμα~\ref{fig:boxplot_load} για το ηλεκτρικό φορτίο, η στατιστική απόκλιση μεταξύ μονοβηματικών και πολυβηματικών στρατηγικών διευρύνεται σημαντικά. Η σταθερότητα και η κυκλικότητα του φορτίου επιτρέπουν στα μονοβηματικά μοντέλα να αξιοποιούν την προηγούμενη τιμή, επιτυγχάνοντας πολύ μικρότερα και λιγότερο ακραία σφάλματα. Οι απολήξεις (whiskers) των καλών μοντέλων συγκρατούνται αναλογικά πολύ κοντά στις διαμέσους τους. Στον αντίποδα, οι πολυβηματικές στρατηγικές (MIMO και Direct) αποτυγχάνουν με κατανομές τεράστιου εύρους και ακραίες απολήξεις που αγγίζουν τα 700~MW.

Η παράλληλη εξέταση των δύο σχημάτων αναδεικνύει δύο χαρακτηριστικά. Πρώτον, η διαφορά απόδοσης μεταξύ μονοβηματικών και πολυβηματικών μοντέλων είναι μεγαλύτερη στο ηλεκτρικό φορτίο από ό,τι στην τιμή DAM.Στο φορτίο η απώλεια της άμεσης αναδρομικής πληροφορίας στις MIMO και Direct διευρύνει το εύρος της κατανομής και των απολήξεων, σε αντίθεση με τη συγκρατημένη εικόνα της τιμής. Δεύτερον, η Recursive μεταβάλλει συμπεριφορά ανάλογα με το σήμα: στην ευμετάβλητη τιμή DAM δυσκολεύεται και πλησιάζει τις πολυβηματικές δομές, ενώ στο σταθερό και κυκλικό φορτίο προσεγγίζει την Teacher-Forcing.

\subsection{Χρονικές και Κυκλικές Συνθήκες Πρόβλεψης}
\label{subsubsect:temporal_hourly_conditions}

Η ωριαία μεροληψία (bias) μετρά τη συστηματική τάση των μοντέλων να υποεκτιμούν ή να υπερεκτιμούν την πραγματική τιμή του σήματος σε κάθε ώρα του 24ώρου. Η ανάλυση αυτή δείχνει την κατεύθυνση του σφάλματος και, κατά συνέπεια, τυχόν δομικά σφάλματα σε συγκεκριμένες περιόδους της ημέρας, ενώ η ερμηνεία της υποεκτίμησης ή υπερεκτίμησης ορίζεται βάσει του προσήμου της διαφοράς.

\begin{figure}[H]
  \centering
  \includegraphics[width=\textwidth]{ch6_forecast_analysis/19_bias_price.png}
  \caption{Προφίλ ωριαίας μεροληψίας (Bias) στην τιμή DAM, Q1 2026. Θετικές τιμές υποδηλώνουν υποεκτίμηση και αρνητικές τιμές υπερεκτίμηση.}
  \label{fig:bias_price}
\end{figure}

Το Σχήμα~\ref{fig:bias_price} δείχνει την ωριαία μεροληψία για την τιμή DAM. Όλες οι στρατηγικές παρουσιάζουν συστηματική υπερεκτίμηση (αρνητικό bias) κατά το μεγαλύτερο μέρος της ημέρας. Το φαινόμενο κλιμακώνεται στις πολυβηματικές προσεγγίσεις MIMO και Direct, οι οποίες αγγίζουν αρνητικές τιμές έως $-25$~€/MWh. Η συμπεριφορά αντιστρέφεται απότομα μόνο στις μεγάλες αιχμές (ώρα 7 και ώρες 19-21), όπου όλα τα μοντέλα εμφανίζουν θετικό bias (υποεκτίμηση), αποτυγχάνοντας να φτάσουν την κορυφή των spikes της αγοράς.

\begin{figure}[H]
  \centering
  \includegraphics[width=\textwidth]{ch6_forecast_analysis/20_bias_load.png}
  \caption{Προφίλ ωριαίας μεροληψίας (Bias) στο ηλεκτρικό φορτίο, Q1 2026. Θετικές τιμές υποδηλώνουν υποεκτίμηση και αρνητικές τιμές υπερεκτίμηση.}
  \label{fig:bias_load}
\end{figure}

Στο Σχήμα~\ref{fig:bias_load} αποτυπώνεται η ωριαία μεροληψία για το ηλεκτρικό φορτίο. Η μεροληψία ακολουθεί την κυκλικότητα της ζήτησης, εναλλάσσοντας πρόσημο μέσα στο 24ωρο. Όλα τα μοντέλα παρουσιάζουν υπερεκτίμηση (αρνητικό bias) νωρίς το πρωί (5-6) και το απόγευμα (16-17), ενώ γυρίζουν σε υποεκτίμηση (θετικό bias) το μεσημέρι (12) και το βράδυ (21). Οι πολυβηματικές στρατηγικές εμφανίζουν εμφανώς μεγαλύτερο εύρος διακύμανσης από τις μονοβηματικές.

Στις ώρες αιχμής, το πρωί (08:00-10:00) και κυρίως το βράδυ (18:00-21:00), η πλειοψηφία των μοντέλων τείνει προς υποεκτίμηση τόσο στην τιμή όσο και στο φορτίο. Αυτό σημαίνει ότι κατά την απότομη άνοδο της ζήτησης και των τιμών, οι στρατηγικές υποτιμούν συστηματικά τόσο το μέγεθος όσο και την ταχύτητα αυτής.

Το ωριαίο προφίλ σφάλματος (Hourly MAE) υπολογίζει τη μέση ακρίβεια των προβλέψεων ανά ώρα, αποκαλύπτοντας αν το σφάλμα ακολουθεί την κυκλικότητα της αγοράς και σε ποια ακριβώς σημεία το μοντέλο αποτυγχάνει συστηματικά. 

\begin{figure}[H]
  \centering
  \includegraphics[width=\textwidth]{ch6_forecast_analysis/07_hourly_price.png}
  \caption{Ωριαίο προφίλ σφάλματος (Hourly MAE) στην τιμή DAM, Q1 2026.}
  \label{fig:hourly_price}
\end{figure}

Όπως φαίνεται στο γράφημα του Σχήματος~\ref{fig:hourly_price}, τα μεγαλύτερα σφάλματα εντοπίζονται στις ώρες αιχμής (πρωί και βράδυ), όπου το MAE ξεπερνά τα $30$~€. Σε αυτές τις συνθήκες υψηλής μεταβλητότητας, οι στρατηγικές Direct, MIMO και Recursive αστοχούν, ενώ η Teacher-Forcing (TF) διατηρεί την καλύτερη επίδοση σε όλο το 24ωρο.

\begin{figure}[H]
  \centering
  \includegraphics[width=\textwidth]{ch6_forecast_analysis/08_hourly_load.png}
  \caption{Ωριαίο προφίλ σφάλματος (Hourly MAE) στο ηλεκτρικό φορτίο, Q1 2026.}
  \label{fig:hourly_load}
\end{figure}

Η ανάλυση των Hourly MAE και Bias αποκαλύπτει ότι η φύση των μεγάλων σφαλμάτων διαφέρει ριζικά ανάμεσα στα δύο προβλήματα. Στην περίπτωση της τιμής, η έξαρση του σφάλματος στις ώρες αιχμής εμφανίζει τη χαρακτηριστική μορφή «Μ», η οποία συμπίπτει με τη συστηματική υποεκτίμηση (under-prediction) των κορυφών. Στο φορτίο, το μεγάλο σφάλμα συνδέεται καθολικά με τη συστηματική υπερεκτίμηση (over-prediction) της ζήτησης κατά τις μεσημεριανές ώρες. Ωστόσο, δεν παρατηρείται κοινή μορφή για όλες τις στρατηγικές, καθώς οι πολυβηματικές προσεγγίσεις (Direct, MIMO) παρουσιάζουν μόνιμη και εντονότερη κατακόρυφη αύξηση της αστοχίας.

\subsection{Ανθεκτικότητα σε Ακραίες Συνθήκες Αιχμής (Spikes)}
\label{subsubsect:spikes_resilience}

Ως spikes τιμής ορίζονται όσες ώρες (22 συνολικά σε όλο το πρώτο τρίμηνο του 2026) η πραγματική τιμή ξεπέρασε το 99ο εκατοστημόριο ($246$~€/MWh). Οι συγκεκριμένες 22 χρονικές στιγμές επισημαίνονται με έντονα κόκκινα σημεία στο Σχήμα~\ref{fig:spikes_timeline_price}.

\begin{figure}[H]
  \centering
  \includegraphics[width=\textwidth]{ch6_forecast_analysis/31a_spikes_timeline_price.png}
  \caption{Χρονική κατανομή και εντοπισμός των ακραίων εξάρσεων (spikes) της τιμής DAM, Q1 2026.}
  \label{fig:spikes_timeline_price}
\end{figure}

Το κριτήριο αστοχίας (miss rate) για την ακρίβεια πρόβλεψης των αιχμών ορίζεται σε απόλυτο σφάλμα άνω των $20$~€/MWh. Όλες οι στρατηγικές καταγράφουν υψηλά ποσοστά αποτυχίας, επιβεβαιώνοντας την αδυναμία των μοντέλων να προβλέψουν με ακρίβεια τις εξάρσεις, όπως αποτυπώνεται στο ραβδόγραμμα του Σχήματος~\ref{fig:spikes_miss_rate_price}.

\begin{figure}[H]
  \centering
  \includegraphics[width=0.72\textwidth]{ch6_forecast_analysis/31b_spikes_miss_rate_price.png}
  \caption{Ποσοστό αστοχίας (Miss Rate) των μοντέλων κατά τις ώρες αιχμής της τιμής DAM.}
  \label{fig:spikes_miss_rate_price}
\end{figure}

Ως spikes φορτίου ορίζονται οι ώρες όπου η ζήτηση ξεπέρασε το 99ο εκατοστημόριο ($8.060$~MW). Προκύπτει συμπτωματικά το ίδιο ακριβώς πλήθος (22 ώρες αιχμής) με την τιμή DAM, επιτρέποντας μια δίκαιη σύγκριση των δειγμάτων. Ωστόσο, οι ώρες αυτές δεν είναι διάσπαρτες, αλλά συγκεντρώνονται όλες μέσα σε ένα χρονικό παράθυρο λίγων μόνο ημερών στα μέσα Ιανουαρίου, λόγω των ακραίων καιρικών συνθηκών που καταπόνησαν το σύστημα, όπως αποτυπώνεται στη χρονοσειρά του Σχήματος~\ref{fig:spikes_timeline_load}.

\begin{figure}[H]
  \centering
  \includegraphics[width=\textwidth]{ch6_forecast_analysis/32a_spikes_timeline_load.png}
  \caption{Χρονική συγκέντρωση των ακραίων εξάρσεων (spikes) του ηλεκτρικού φορτίου, Q1 2026.}
  \label{fig:spikes_timeline_load}
\end{figure}

Στο ραβδόγραμμα του Σχήματος~\ref{fig:spikes_miss_rate_load}, η εικόνα ανατρέπεται ριζικά. Η Teacher-Forcing (WF-TF) σημειώνει τη χαμηλότερη αστοχία με miss rate μόλις $5\%$, η Recursive ακολουθεί με $14\%$, ενώ οι MIMO ($64\%$) και Direct ($100\%$) αποτυγχάνουν καθολικά.

\begin{figure}[H]
  \centering
  \includegraphics[width=0.72\textwidth]{ch6_forecast_analysis/32b_spikes_miss_rate_load.png}
  \caption{Ποσοστό αστοχίας (Miss Rate) των μοντέλων κατά τις ώρες ακραίου ηλεκτρικού φορτίου.}
  \label{fig:spikes_miss_rate_load}
\end{figure}

Η σύγκριση παραμένει δίκαιη: τα $200$~MW αντιπροσωπεύουν μόλις το $2,5\%$ της ακραίας ζήτησης, έναντι του $8\%$ της τιμής ($20$~€/MWh). Παρότι ο έλεγχος είναι αυστηρότερος στο φορτίο, αυτό εξακολουθεί να παρουσιάζει πολύ χαμηλότερη αστοχία, επιβεβαιώνοντας ότι είναι πιο προβλέψιμο μέγεθος από την τιμή.

Η παραπάνω στατιστική συμπεριφορά γίνεται άμεσα κατανοητή μέσω της ποιοτικής εξέτασης δύο χαρακτηριστικών επεισοδίων στα Σχήματα~\ref{fig:spike_episode_price} (τιμή DAM) και~\ref{fig:spike_episode_load} (ηλεκτρικό φορτίο), τα οποία αποτυπώνουν τον τρόπο με τον οποίο οι στρατηγικές αντιδρούν σε πραγματικό χρόνο σε επεισόδια αιχμής.

\begin{figure}[H]
  \centering
  \includegraphics[width=\textwidth]{ch6_forecast_analysis/31c_spike_episode_price.png}
  \caption{Πραγματική απόκριση των στρατηγικών πρόβλεψης κατά τη διάρκεια ενός επεισοδίου έξαρσης (spike) της τιμής DAM.}
  \label{fig:spike_episode_price}
\end{figure}

Στην τιμή DAM (Σχήμα~\ref{fig:spike_episode_price}), κατά την κορυφή της έξαρσης ($>270$~€/MWh), η Teacher-Forcing (WF-TF) είναι η μόνη που ακολουθεί την πορεία της ανόδου, περιορίζοντας τη μονοβηματική αστοχία στην ελάχιστη υποεκτίμηση των $-42$~€/MWh. Αντίθετα, οι υπόλοιπες στρατηγικές αποτυγχάνουν πλήρως να ακολουθήσουν το spike, με τη Recursive να μένει χαμηλά και τις πολυβηματικές (MIMO, Direct) να καταγράφουν σφάλματα που ξεπερνούν κατά πολύ τα $100$~€/MWh.
\begin{figure}[H]
  \centering
  \includegraphics[width=\textwidth]{ch6_forecast_analysis/32d_spike_episode_load.png}
  \caption{Πραγματική απόκριση των στρατηγικών πρόβλεψης κατά τη διάρκεια ενός επεισοδίου αιχμής του ηλεκτρικού φορτίου.}
  \label{fig:spike_episode_load}
\end{figure}


Στο ηλεκτρικό φορτίο (Σχήμα~\ref{fig:spike_episode_load}), όπου η πραγματική αιχμή προσεγγίζει τα $9.000$~MW, η Teacher-Forcing και η Recursive ακολουθούν ικανοποιητικά την καμπύλη, με την καλύτερη προσέγγιση να σημειώνει ελάχιστη υποεκτίμηση $-186$~MW. Στον αντίποδα, οι πολυβηματικές MIMO και Direct αδυνατούν να ακολουθήσουν την κορυφή, εμφανίζοντας έντονο flat-lining που ξεπερνά το όριο αστοχίας των $200$~MW.

Συνολικά, από τα παραπάνω συμπεραίνουμε ότι οι μονοβηματικές στρατηγικές προσεγγίζουν καλύτερα τις ακραίες συνθήκες (spikes) καθώς διαθέτουν πιο πρόσφατη πληροφορία, ενώ οι πολυβηματικές κρίνονται ακατάλληλες να διαχειριστούν τέτοιες ακραίες καταστάσεις και, κατ' επέκταση, το ρίσκο, τόσο όσον αφορά την τιμή αλλά και το φορτίο, παρόλη την περιοδικότητα του.





\section{Στατιστική Επικύρωση των Προβλέψεων (Diebold-Mariano)}
\label{sec:dm_validation}

Οι κατατάξεις των προηγούμενων ενοτήτων εγείρουν ερώτημα εγκυρότητας: Αντανακλούν οι 
διαφορές σφάλματος ουσιαστική υπεροχή ή θόρυβο της δοκιμαστικής περιόδου; Ο έλεγχος
Diebold-Mariano (DM) απαντά συγκρίνοντας κατά ζεύγη τα σφάλματα δύο στρατηγικών και
αποδίδοντας στατιστική σημασία στη διαφορά τους. Εφαρμόζουμε τη διόρθωση
Harvey-Leybourne-Newbold για μικρά δείγματα και τη διόρθωση Holm-Bonferroni για τις
πολλαπλές συγκρίσεις, ώστε κάθε «νίκη» να στηρίζεται σε αυστηρό στατιστικό θεμέλιο
($T = 2.160$ ώρες, $p = 0{,}05$).

\begin{figure}[H]
  \centering
  \includegraphics[width=\textwidth]{ch5_dm_test/dm_price_wins_v2.png}
  \caption{Κατάταξη μοντέλων βάσει στατιστικά σημαντικών νικών DM-HLN ($p<0{,}05$)
           έναντι των υπολοίπων, τιμή DAM, Q1 2026.}
  \label{fig:dm_price_wins_c6}
\end{figure}

Στην τιμή DAM (Σχ.~\ref{fig:dm_price_wins_c6}) η Walk-Forward TF υπερτερεί στατιστικά
και των 18 υπολοίπων μοντέλων ($p < 0{,}001$). Εντός της ομάδας Walk-Forward TF
οι αλγορίθμοι δεν απέχουν στατιστικά μεταξύ τους, κάτι που επιβεβαιώνει ξανά
ότι η επιλογή στρατηγικής είναι διαχωριστικός παράγοντας.
Ακολουθούν η στατική TF, οι αναδρομικές και τελευταίες οι MIMO/Direct, που δεν
ξεπερνούν στατιστικά καμία μονοβηματική προσέγγιση. Ο έλεγχος επικυρώνει την κατάταξη
της §\ref{sec:strategic_evaluation}: η υπεροχή της walk-forward επανεκπαίδευσης δεν
οφείλεται στην τύχη της περιόδου.

\begin{figure}[H]
  \centering
  \includegraphics[width=\textwidth]{ch5_dm_test/dm_load_wins_v2.png}
  \caption{Κατάταξη μοντέλων βάσει στατιστικά σημαντικών νικών DM-HLN ($p<0{,}05$)
           έναντι των υπολοίπων, ηλεκτρικό φορτίο, Q1 2026.}
  \label{fig:dm_load_wins_c6}
\end{figure}

Στο φορτίο (Σχ.~\ref{fig:dm_load_wins_c6}) η εικόνα αλλάζει στο κρίσιμο σημείο: η
Walk-Forward αναδρομική (Rec-XGB-SS) είναι στατιστικά αδιάκριτη από την ιδεατή
Walk-Forward TF, παρότι δεν διαθέτει πρόσβαση στις πραγματικές
ενδοημερήσιες τιμές. Η ισχυρή ενδοημερήσια αυτοσυσχέτιση του φορτίου συγκρατεί τη συσσώρευση σφάλματος, ενώ στην τιμή οι αιχμές τη μεγεθύνουν. Οι MIMO και
Direct παραμένουν στατιστικά κατώτερες κάθε μονοβηματικής στρατηγικής.

Συνολικά, ο έλεγχος DM μετατρέπει τις αριθμητικές διαφορές των πινάκων σε στατιστικά
τεκμηριωμένα συμπεράσματα: η υπεροχή της walk-forward επανεκπαίδευσης στην τιμή και η
ισοδυναμία αναδρομικής και Teacher-Forcing στο φορτίο είναι ανθεκτικές ιδιότητες, όχι
τυχαία ευρήματα της δοκιμαστικής περιόδου.

% ---------------------------------------------------------------
\section{Σύγκριση με Σχετικές Εργασίες}
\label{sec:related_comparison}

Αντιπαραβάλλουμε τα ευρήματά μας με τη διεθνή και την εγχώρια βιβλιογραφία, εστιάζοντας
 στα σημεία σύγκλισης και απόκλισης με μελέτες πρόβλεψης τιμής
και φορτίου με έμφαση στην ελληνική αγορά. Tο μεθοδολογικό πλαίσιο των εργασιών αυτών συνοψίζεται
στον Πίνακα~\ref{tab:related_comparison} του Κεφαλαίου~\ref{ch:literature}.

\subsection{Μεθοδολογική ευθυγράμμιση με την ελληνική αγορά}
\label{subsec:greek_alignment}

Οι εγχώριες μελέτες εξετάζουν κάθε σήμα μεμονωμένα: ο Kanousis~\cite{kanousis2025}
και ο Κοντράρος~\cite{kontraros2023} την τιμή της Αγοράς Επόμενης Ημέρας, ο
Μαγκλαράς~\cite{magklaras2018} καθώς και οι Shiblee και Koukaras~\cite{shiblee2025shortterm} 
το ηλεκτρικό φορτίο. Η παρούσα εργασία μοντελοποιεί ταυτόχρονα και τα δύο σήματα 
και αξιοποιεί τη μεταξύ τους σχέση, τροφοδοτώντας την προβλεπόμενη ζήτηση ως είσοδο για την τιμή. 

Η δική μας αξιολόγηση έδειξε ότι η αντικατάσταση της ακριβέστερης πρόβλεψης (MAE = 80,7~MW) 
με την ENTSO‐E πρόβλεψη (MAE = 267~MW) μεταβάλλει το τελικό MAE τιμής κατά < $0,1\%$, 
υποδηλώνοντας ότι τα μοντέλα τιμής βασίζονται κυρίως στις χρονικές υστερήσεις τιμής. 
Το εύρημα αυτό, ότι οι προσεγγίσεις Μηχανικής Μάθησης και συγκεκριμένα οι δομές Gradient 
Boosting (LightGBM) μπορούν να ξεπεράσουν αισθητά τις επίσημες επιχειρησιακές προβλέψεις 
της πλατφόρμας ENTSO-E στο ελληνικό σύστημα, επιβεβαιώνεται καθολικά και από τους 
Shiblee και Koukaras~\cite{shiblee2025shortterm}, οι οποίοι ανέδειξαν επίσης το αυτόνομο 
LightGBM ως το κορυφαίο μοντέλο για το φορτίο.

Σε επίπεδο τάξης μεγέθους, τα σφάλματα της παρούσας εργασίας κινούνται στο ίδιο 
εύρος με τις εγχώριες μελέτες της ελληνικής αγοράς και του φορτίου. Ενδεικτικά, οι τιμές 
σφαλμάτων για την τιμή DAM (όπου στη δική μας αναδρομική LightGBM υλοποίηση καταγράφεται 
MAE 14,028~€/MWh, έναντι MAE 12,502~€/MWh του Kanousis~\cite{kanousis2025} και 
MAE 12,030~€/MWh της μελέτης του ΕΜΠ~\cite{ntua2025shortwindow}) και για το ηλεκτρικό 
φορτίο (όπου οι Shiblee και Koukaras σημειώνουν test MAE 69,12~MW και η δική μας 
Recursive LightGBM υλοποίηση MAE 80,704~MW) κινούνται σε εξίσου χαμηλά επίπεδα. 
Η διαφορετική περίοδος δοκιμής, το επίπεδο και η μεταβλητότητα των τιμών/φορτίου, 
καθώς και ο ορισμός του σφάλματος καθιστούν κάθε αριθμητική αντιπαραβολή ενδεικτική 
και όχι αυστηρή (Πίν.~\ref{tab:related_comparison}).

Η ισχυρότερη σύγκλιση προκύπτει με τη μελέτη της ομάδας του ΕΜΠ~\cite{ntua2025shortwindow},
η οποία αξιολογεί LSTM, XGBoost, LightGBM και CatBoost σε Ελλάδα, Βέλγιο και Ιρλανδία
υπό το ίδιο καθεστώς προβλεπόμενων χαρακτηριστικών. Τα ευρήματά μας επιβεβαιώνουν δύο
πορίσματά της: το LightGBM παραμένει το ακριβέστερο και αποδοτικότερο μοντέλο, και η χρονική
εγγύτητα των δεδομένων εκπαίδευσης βελτιώνει την προσαρμογή στις εποχικές μεταβολές
και τις αιχμές, όπως ποσοτικοποιεί το Retraining Premium μας (Σχ.~\ref{fig:delta_price}).
Η διαφορά εντοπίζεται στον μηχανισμό ανανέωσης, όχι στο συμπέρασμα: οι ερευνητές του ΕΜΠ επιτυγχάνουν
στηρίζονται αποκλειστίκα σε πρόσφατα μέσω σύντομων κυλιόμενων παραθύρων με "νικητή" το παράθυρο 45-60 ημερών
ενώ εμείς διατηρούμε ολόκληρο το ιστορικό και το εμπλουτίζουμε μηνιαία. Παρά τη διαφορά, το συμπέρασμα ταυτίζεται: στη μεταβλητή
ελληνική αγορά η ανανέωση των δεδομένων έχει μεγαλύτερη επιρροή από κάθε άλλη σχεδιαστική
επιλογή.

\subsection{Αναδρομική πρόβλεψη και η θεωρία συσσώρευσης σφάλματος}
\label{subsec:rec_vs_taieb}

Η κλασική ταξινόμηση των στρατηγικών~\cite{taieb2012review,bontempi2012strategies}
αποδίδει στην αναδρομική πρόβλεψη δομική μειονεξία έναντι των multiple output
αρχιτεκτονικών, λόγω της συσσώρευσης σφάλματος κατά την
ανατροφοδότηση~\cite{bento2024mimo}. Τα ευρήματά μας για το φορτίο αντιστρέφουν τη
θέση αυτή: η αναδρομική walk-forward κυριαρχεί των MIMO και Direct και πλησιάζει το
ιδεατό άνω φράγμα (Πίνακας~\ref{tab:wf_tf_load}), ενώ οι πολυέξοδες δομές
εκτινάσσονται ακριβώς στις ημερήσιες αιχμές (Σχ.~\ref{fig:hourly_load}), όπου λείπει
η ενδοημερήσια ανατροφοδότηση. 

Nα σημειωθεί ότι στην παρούσα εργασία η Recursive στρατηγική χωρίζεται σε αυτόνομα 24ωρα. Με το καθημερινό reset στο gate closure, η αναδρομική διαδικασία περιορίζεται αυστηρά εντός της ημέρας κοιτάζοντας το πολύ έως 23 ώρες πίσω, οπότε το σφάλμα δεν προλαβαίνει να συσσωρευτεί σε μεγάλο βαθμό.

\subsection{Επανεκπαίδευση και μεταβολή καθεστώτος}
\label{subsec:retrain_regime}

Η βιβλιογραφία αναγνωρίζει ποιοτικά και ποσοτικά ότι τα στατικά μοντέλα χάνουν ακρίβεια στις εποχικές μεταβάσεις. Στην εργασία μας, η αξία της μηνιαίας επανεκπαίδευσης αποτυπώνεται αρχικά στο Σχήμα~\ref{fig:monthly_price}, όπου η walk-forward προσέγγιση συγκρατεί το σφάλμα της τιμής DAM σε χαμηλότερα επίπεδα σε σχέση με τη στατική εκπαίδευση, με τη διαφοροποίηση να γίνεται πιο ορατή από τον Ιανουάριο.

Η συνολική αποτίμηση αυτού του οφέλους ανά στρατηγική επιβεβαιώνεται από τα διαγράμματα διαφοράς επίδοσης (Σχήματα~\ref{fig:delta_price} και~\ref{fig:delta_load}). Τα αποτελέσματα δείχνουν ότι η επανεκπαίδευση έχει ουσιαστικό αντίκτυπο στην τιμή DAM. Αντίθετα, στο ηλεκτρικό φορτίο η επίδρασή της είναι μικρότερη διότι, όπως εξηγήσαμε βάσει του Σχήματος~\ref{fig:regime_load}, δεν υπήρχε μεταβολή καθεστώτος.

Το εύρημα αυτό συμφωνεί με τις παρατηρήσεις της ομάδας του ΕΜΠ~\cite{ntua2025shortwindow} για την ανάγκη χρήσης πρόσφατων δεδομένων στην ελληνική αγορά, επεκτείνοντάς το σε τέσσερις διαφορετικές στρατηγικές πρόβλεψης.

\subsection{Πολυπλοκότητα μοντέλου και αξία των συνδυαστικών μετρικών}
\label{subsec:complexity_metrics}

Έναντι προσεγγίσεων που στρέφονται σε σύνθετα δίκτυα βαθιάς μάθησης, τα ευρήματά μας
δείχνουν ότι ελαφριά tree-based μοντέλα, με κατάλληλη αρχιτεκτονική εισόδων (πυκνές
υστερήσεις) και walk-forward επανεκπαίδευση, επιτυγχάνουν εφάμιλλη ή ανώτερη ακρίβεια.
Η μηχανική χαρακτηριστικών υπερτερεί της αλγοριθμικής πολυπλοκότητας, συμβατά με την
έμφαση του Κοντράρου~\cite{kontraros2023} στην επιλογή χαρακτηριστικών και με την
πρακτική των ανοιχτών σημείων αναφοράς~\cite{lago2021}.

% ---------------------------------------------------------------
\section{Σύνοψη Αποτελεσμάτων και Συμπεράσματα}
\label{sec:ch6_summary}

Η αξιολόγηση του κεφαλαίου οδηγεί σε ένα κεντρικό συμπέρασμα: το χρονικό πλαίσιο
εκπαίδευσης και η μηχανική χαρακτηριστικών έχουν μεγαλύτερη βαρύτητα από την επιλογή
αλγορίθμου. Η μηνιαία walk-forward επανεκπαίδευση μειώνει το MAE της τιμής DAM κατά
27,7\% (14,478~$\rightarrow$~10,463~€/MWh), με ανάλογη πτώση σε RMSE και sMAPE, ενώ στο
φορτίο το όφελος είναι προς αμφισβήτηση. Η ασυμμετρία εξηγείται από τη δομική μετατόπιση του
επιπέδου της τιμής τον Φεβρουάριο 2026 (§\ref{subsec:regime_evidence}), την οποία το
ομαλό φορτίο δεν εμφανίζει· σε κοινή ποσοστιαία κλίμακα η επανεκπαίδευση κερδίζει
$-4{,}496$ μονάδες sMAPE στην τιμή έναντι $-0{,}529$ στο φορτίο.

Ανάμεσα στις εφαρμόσιμες για πρόβλεψη DAM στρατηγικές, η αναδρομική (Recursive) είναι η
ακριβέστερη (Σχ.~\ref{fig:dm_price_wins_c6}). Ο ημερήσιος μηδενισμός της αναδρομικής
διαδικασίας περιορίζει τη συσσώρευση σφάλματος εντός του 24ώρου, αναιρώντας στην πράξη
το θεωρητικό μειονέκτημα που της αποδίδει η κλασική ταξινόμηση έναντι των multistep
δομών. Στο φορτίο, όπου η ενδοημερήσια αυτοσυσχέτιση είναι ισχυρή, η Recursive
walk-forward πλησιάζει το ιδεατό Teacher-Forcing (70,641 έναντι 67,729~MW) και ο έλεγχος
Diebold-Mariano τις βρίσκει στατιστικά αδιάκριτες ($p=0{,}769$,
§\ref{sec:dm_validation}, Σχ.~\ref{fig:dm_load_wins_c6}).

Στις ακραίες συνθήκες η διαφορά γίνεται καθοριστική. Στις αιχμές ζήτησης (spikes) του
Ιανουαρίου, η Direct αστοχεί στο 100\% των ωρών αιχμής και η MIMO στο 64\%, έναντι 14\%
της Recursive και 5\% της Teacher-Forcing (Σχ.~\ref{fig:spikes_miss_rate_load}).
Η αστοχία τους οφείλεται σε συστηματική υποεκτίμηση συγκεντρωμένη στις
ώρες αιχμής (Σχ.~\ref{fig:hourly_price},~\ref{fig:bias_price}), όπου λείπει η ενδοημερήσια
ανατροφοδότηση. Το ηλεκτρικό φορτίο παραμένει συνολικά πιο προβλέψιμο από την τιμή
(sMAPE 1,3\% έναντι 10,7\%), ακόμη και υπό αυστηρότερο κριτήριο αιχμής.

Ο έλεγχος Diebold-Mariano (§\ref{sec:dm_validation}) επικυρώνει ότι οι διαφορές αυτές
είναι ανθεκτικές και όχι τυχαίες. Το Κεφάλαιο~\ref{ch:fi} ερμηνεύει το γιατί μέσω της
ανάλυσης σημαντικότητας χαρακτηριστικών, αναδεικνύοντας ποια χαρακτηριστικά και σε τι
ποσοστό αξιοποιεί κάθε στρατηγική.
